% Week 12 实验周报 —— 编译方式：XeLaTeX（数字来自 docs/week12/numbers.tex）
\documentclass[11pt]{ctexart}
\usepackage[a4paper, margin=2.5cm]{geometry}
\usepackage{booktabs}
\usepackage{hyperref}
\usepackage{enumitem}
\setlist{nosep}
% generated by analysis/week12/04_report_numbers.py - do not edit by hand
\newcommand{\rnWnineQuoteMs}{1.8}
\newcommand{\rnWnineRedisMs}{0.15}
\newcommand{\rnWnineRenderMs}{724}
\newcommand{\rnWnineDailyS}{4.5}
\newcommand{\rnMemIntervals}{670}
\newcommand{\rnMemTickers}{666}
\newcommand{\rnPanelTenRows}{617{,}211}
\newcommand{\rnPanelTenTickers}{532}
\newcommand{\rnEone}{9.5}
\newcommand{\rnEtwo}{6.8}
\newcommand{\rnEthree}{5.4}
\newcommand{\rnEfour}{163.8}
\newcommand{\rnEfive}{69.9}
\newcommand{\rnTenDailyS}{38.0}
\newcommand{\rnGoneDWin}{72}
\newcommand{\rnGoneDPooled}{0.812}
\newcommand{\rnRoundOneRidge}{7.16}
\newcommand{\rnRoundOneGbdt}{2.23}
\newcommand{\rnNoPrices}{51}
\newcommand{\rnInStudy}{597}
\newcommand{\rnTruncated}{3}
\newcommand{\rnPostExit}{7}
\newcommand{\rnLatestBeforeMs}{457}
\newcommand{\rnLatestBeforeRows}{579}
\newcommand{\rnLatestFinalMs}{0.3}
\newcommand{\rnLatestFinalRows}{1496}
\newcommand{\rnQuBeforeMs}{390}
\newcommand{\rnQuBeforeN}{578}
\newcommand{\rnQuAfterMs}{23.5}
\newcommand{\rnQuAfterN}{1496}
\newcommand{\rnMarketBeforeS}{12.8}
\newcommand{\rnMarketAfterS}{2.6}
\newcommand{\rnFeatBeforeS}{14.7}
\newcommand{\rnFeatAfterS}{14.5}
\newcommand{\rnDailyAfterS}{25.8}
\newcommand{\rnScaleSteadyS}{59.5}
\newcommand{\rnScaleN}{1517}
\newcommand{\rnScaleBackfillS}{530}
\newcommand{\rnRenderBeforeMs}{591}
\newcommand{\rnRenderAfterMs}{514}
\newcommand{\rnFinalQuMs}{50}
\newcommand{\rnFinalQuN}{1496}
\newcommand{\rnFinalScalarMs}{27}
\newcommand{\rnFinalXsMs}{27}
\newcommand{\rnFinalRenderMs}{510}
\newcommand{\rnGarchOrdS}{3.2}
\newcommand{\rnGarchRefitS}{15.2}
\newcommand{\rnGarchRefits}{1{,}326}
\newcommand{\rnGarchTickers}{1{,}479}
\newcommand{\rnCatchupS}{84.4}
\newcommand{\rnCatchupFeatS}{62.7}
\newcommand{\rnDbFinalMb}{1601}
\newcommand{\rnDbElevenMb}{1570}
\newcommand{\rnRssFinalMb}{299}
\newcommand{\rnRssElevenMb}{286}
\newcommand{\rnElevenTestMl}{0.755}
\newcommand{\rnElevenTestG}{0.791}
\newcommand{\rnElevenMaeMl}{7.91}
\newcommand{\rnElevenMaeG}{8.36}
\newcommand{\rnElevenMaeE}{10.35}
\newcommand{\rnElevenVal}{0.705}
\newcommand{\rnElevenTestStart}{2024-05-24}
\newcommand{\rnNoDummyTest}{0.741}
\newcommand{\rnUnknownN}{52}
\newcommand{\rnUnknownRemoved}{51}
\newcommand{\rnGTrain}{0.840}
\newcommand{\rnGVal}{0.777}
\newcommand{\rnCsOnly}{0.742}
\newcommand{\rnLevelOnly}{0.801}
\newcommand{\rnOosStart}{2022-12-27}
\newcommand{\rnOosEnd}{2024-12-31}
\newcommand{\rnOosDates}{506}
\newcommand{\rnOosRows}{244{,}038}
\newcommand{\rnOosTickers}{505}
\newcommand{\rnOosMaeE}{7.55}
\newcommand{\rnOosMaeG}{6.49}
\newcommand{\rnOosMaeMl}{6.15}
\newcommand{\rnHOneDelta}{1.07}
\newcommand{\rnHOneLo}{0.65}
\newcommand{\rnHOneHi}{1.62}
\newcommand{\rnHOneSizeTicker}{48}
\newcommand{\rnHOneSizeSector}{9.3}
\newcommand{\rnHOnePholm}{0.002}
\newcommand{\rnHTwoDelta}{0.337}
\newcommand{\rnHTwoLo}{0.206}
\newcommand{\rnHTwoHi}{0.456}
\newcommand{\rnHTwoSizeTicker}{50}
\newcommand{\rnHTwoSizeSector}{4.7}
\newcommand{\rnHTwoPholm}{0.002}
\newcommand{\rnHThreeDelta}{1.48}
\newcommand{\rnHThreeLo}{1.15}
\newcommand{\rnHThreeHi}{1.80}
\newcommand{\rnHThreeSizeTicker}{23}
\newcommand{\rnHThreeSizeSector}{3.7}
\newcommand{\rnHThreePholm}{0.002}
\newcommand{\rnHFourDelta}{1.77}
\newcommand{\rnHFourLo}{1.29}
\newcommand{\rnHFourHi}{2.44}
\newcommand{\rnHFourSizeTicker}{42}
\newcommand{\rnHFourSizeSector}{8.0}
\newcommand{\rnHFourPholm}{0.002}
\newcommand{\rnHTwoYThree}{0.048}
\newcommand{\rnHTwoYFour}{0.634}
\newcommand{\rnHTwoBDelta}{0.286}
\newcommand{\rnSelfIidTicker}{6.7}
\newcommand{\rnSelfShockTicker}{69}
\newcommand{\rnSelfShockSector}{9.3}
\newcommand{\rnIvJpm}{0.710}
\newcommand{\rnIvJpmCorr}{0.78}
\newcommand{\rnIvJpmCorrConst}{0.83}
\newcommand{\rnIvPanelRows}{467{,}189}
\newcommand{\rnIvPanelTickers}{506}
\newcommand{\rnIvMaeIv}{5.16}
\newcommand{\rnIvMaeGd}{6.63}
\newcommand{\rnIvMaeGy}{6.92}
\newcommand{\rnIvMaeRv}{7.60}
\newcommand{\rnIvMaeEw}{7.52}
\newcommand{\rnIvWinGd}{92}
\newcommand{\rnIvWinGy}{70}
\newcommand{\rnIvConstIv}{4.69}
\newcommand{\rnIvConstGy}{6.47}
\newcommand{\rnIvXsIv}{0.95}
\newcommand{\rnIvXsGy}{0.84}
\newcommand{\rnIvTsIv}{0.11}
\newcommand{\rnIvGapActual}{1.77}
\newcommand{\rnIvGapConst}{1.78}
\newcommand{\rnBandGCalm}{5.48}
\newcommand{\rnBandGFear}{5.86}
\newcommand{\rnBandESameCalm}{6.66}
\newcommand{\rnBandESameFear}{4.89}
\newcommand{\rnBandEFullCalm}{9.01}
\newcommand{\rnBandEFullFear}{10.83}
\newcommand{\rnBandGDays}{1011}
\newcommand{\rnToolAsOf}{2026-10-02}
\newcommand{\rnToolQuoted}{1{,}493}
\newcommand{\rnToolGarch}{1{,}478}
\newcommand{\rnToolEwma}{15}
\newcommand{\rnToolMl}{496}


\title{\textbf{Chooser Option 定价项目 —— 第 12 周实验报告}\\
\large Enhancement 第 4 周（结项）：回答导师两问、换锚上线、ML 参考价扩展、终测与结项报告}
\author{佟铸城（Zhucheng Tong）}
\date{2026 年 10 月 13 日}

\begin{document}
\maketitle

\section{本周目标}
回答导师对第 11 周的两个问题（ML 增量的统计显著性边界；隐含波动率是否值得作为候选锚），把第 11 周决定的换锚落到工具上，复测全部审计与性能，并在第 8 周最终报告原文件上续写第 9--12 周，形成结项报告。

\section{完成工作}
\begin{enumerate}
  \item \textbf{预注册}：检验判据在跑数前写入 \texttt{docs/week12/preregistration.json} 并记录哈希（假设、统计量、推断程序、size 检验门槛、Holm 矫正、决策规则）；
  \item \textbf{显著性与隐含波动率}：\texttt{src/analysis\_w12.py}；S\&P 500 面板的 OptionMetrics 隐含波动率按收益相关匹配证券代码后提取（仅本机）；
  \item \textbf{换锚上线}：GARCH(1,1) 1 日锚逐股票存储状态、增量续算，主定价切换（新股退回 EWMA），误差带按新锚重算；
  \item \textbf{ML 参考价扩展}：按事先规则，把横截面 GBDT 固化为参考模型，覆盖当日点时成分股；
  \item \textbf{修复}：调度器漏跑、写库清空未提供的列、批量重建无逐票容错、研究面板与新增列冲突、代码被复用的退出成员；
  \item \textbf{终测与审计}：第 12 周 12 项审计全过（连带回归第 7、9--11 周）；性能终测；
  \item \textbf{结项报告}：17 页（\texttt{docs/week12/final\_report\_cn.pdf}），第 7--12 节数字全部由脚本生成为宏。
\end{enumerate}

\section{关键结果}
\begin{itemize}
  \item \textbf{ML 增量的显著性边界}：2023--2024 年样本外，相对 GARCH 锚的增量为每份 \rnHTwoDelta{} 美元 MAE，95\% 区间 \rnHTwoLo--\rnHTwoHi（行业簇 bootstrap，size \rnHTwoSizeSector\%，Holm 后显著）。区间以这段时期为条件；增量主要来自 2024 年（\rnHTwoYFour，2023 年 \rnHTwoYThree），独立年份只有两个，不做时间维度的显著性推断。按事先规则判定成立，ML 参考价扩展到面板成员。
  \item \textbf{隐含波动率的定位}：面板上是检验过的最强的锚（MAE \rnIvMaeIv{} 美元，GARCH 1 日 \rnIvMaeGd、GARCH 1 年 \rnIvMaeGy），优势来自每只股票的波动率水平定位而非择时；对 GARCH 1 日锚的优势显著，对免费平稳锚的优势（\rnHFourDelta{} 美元）的推断程序 size 不合格、只作描述，可部署性无法核实——按事先判据定位为\textbf{研究基准}，不列为候选锚。
  \item \textbf{换锚}：GARCH 1 日锚相对 EWMA 的增量 \rnHOneDelta{} 美元、每个合格年份为正，但推断程序 size 不合格，只作描述；换锚决定依据第 11 周三个窗口的证据。工具误差带改为 \rnBandGCalm{} / \rnBandGFear{} 美元，校准窗口不含 2020 年危机，已注明可能偏窄。
  \item \textbf{性能}：数据库端最新状态读取 \rnLatestBeforeMs{} $\to$ \rnLatestFinalMs{} ms；GARCH 步骤每日成本普通日 \rnGarchOrdS{} s、重估日 \rnGarchRefitS{} s；新 ML 参考价需要当日整个截面，单票报价（未命中缓存）\rnWnineQuoteMs{} $\to$ \rnFinalScalarMs{} ms。
\end{itemize}

\section{实验记录与发现}
\begin{itemize}
  \item size 检验本身先在已知答案的数据上验证，发现第一版有偏（纯噪声下拒绝率明显偏高：行业日均值混入了个股噪声），修正后才用于真实数据；四项假设上按股票重抽样的程序全部严重失准（\rnHTwoSizeTicker\% 等），只有改用行业簇程序且逐项通过 size 检验的才给显著性。
  \item 复核发现第 10 周把“历史不足”归因于新上市或改名有误，实为退市（生存偏差）与代码被他司复用；后者 \rnPostExit{} 个已在工具中停用，研究面板不受影响。
  \item 调度器漏跑的真正原因：宿主机睡眠时虚拟机挂起、单调时钟停走，长等待被不断推迟；改为按墙钟的 10 分钟心跳补跑。
  \item GARCH 参数靠数值优化求得，macOS 与 Linux 容器的结果在末几位不同；“逐位一致”只在同一平台内成立，跨平台审计改用相对容差。
\end{itemize}

\section{交付与后续}
交付：代码包、结项报告、本周报、学习报告（终版）。演示形式（视频或现场）待与导师确认。后续方向（结项报告第 13 节）：用隐含波动率做 GARCH 的方差目标化；更长的样本外历史；美式 chooser 与 ML 加速器。

\end{document}
